% Withdrawal, design implications, and conclusion for "The Unchosen Adjacency"
% Ofcom figures checked against Ofcom's 2026 report and its published summary.

\section{Withdrawal}\label{sec:withdrawal}

If some actions are defective, users who cannot disaggregate the meanings attached to them have three options: accept the surplus, act less, or leave. The last two are forms of withdrawal. This section states a hypothesis about withdrawal and is careful to separate what is observed from what is proposed.

\subsection{Three levels}

\paragraph{Observed decline in contribution.} Ofcom's \emph{Adults' Media Use and Attitudes} report, published 2 April 2026 \cite{ofcom2026}, finds that 49 percent of UK adult social media users now post, share, or comment, down from 61 percent in 2024. Ofcom links the decline to video-focused features and to concern about the long-term consequences of past posts, and the share of adults worried that their posts could cause them problems later rose from 43 to 49 percent. This supports the claim that the dependent variable is real and measurable. It does not identify a cause, and the explanations Ofcom offers are competitors to the bundling mechanism proposed here and not instances of it. Ofcom's measure is also the share of users who post at all, whereas the hypothesis below concerns how often a user performs a given action, so the data bear on the first margin and not the second.

\paragraph{A proposed mechanism.} The hypothesis is that part of the decline is produced by bundling. Where an action carries consequential surplus that the user cannot separate, the user's expected cost of performing it rises, and the rate of performing it falls.

\begin{quote}
\textbf{Hypothesis (withdrawal).} As the share of a user's intended actions that are defective rises, and as attribution $J_o$ for their consequential surplus rises, the rate at which the user performs those actions falls.
\end{quote}

\paragraph{The untested connection.} No evidence presented here links the mechanism to the decline. Competing explanations are plausible and should be listed: a shift from posting to passive consumption, algorithmic ranking that reduces the reach of ordinary posts, privacy and permanence concerns unrelated to bundling, fatigue, and migration to messaging. The hypothesis is one candidate among these.

\subsection{What would test it}

The hypothesis makes differential predictions that the competing explanations do not all share.

\begin{enumerate}
  \item \textbf{Differential decline by action type.} Actions with more consequential mismatch, such as posting where an attached advertisement is likely to be read as the author's, should decline more than actions with less.
  \item \textbf{Response to new controls.} When a platform introduces an action that separates a previously bundled relation, contribution of the affected action should rise relative to comparable unaffected actions. Such introductions are natural experiments.
  \item \textbf{Stated reasons.} Users who reduce an action should report the unwanted association or lost relation as a reason more often than users who do not, controlling for age and platform. One kind of stated reason to record is an objection to the platform-selected content that follows a user's Reels.
\end{enumerate}

Failure of these predictions would count against the hypothesis. A uniform decline unrelated to bundling would favor the competing explanations.

\subsection{The commercial dimension}

If the hypothesis holds, a platform that bundles advertising into user posts draws short-term inventory from the context it appropriates and gradually reduces the willingness to produce the content that makes the inventory valuable. This is a claim about cost to the platform and is conditional on the hypothesis. It is stated to indicate why the bundling may not serve the platform's own interest, not to assert that it does not.

\section{Design Implications}\label{sec:design}

The framework classifies bundles and does not by itself mandate removing them. What follows are candidate separations. Some are supported by comparators discussed in the earlier sections and others are not yet shown to be feasible. Their user-side benefit is described, and their platform-side cost is not measured here.

\begin{enumerate}
  \item \textbf{Transparent participation exclusion that preserves view, on personal profiles.} An action that bars commenting, replying, tagging, and messaging by a chosen person while leaving public material visible, and that carries a stated reason and a duration. Facebook Groups already offer time-limited suspension \cite{fb-groups}, Pages offer bans that leave the Page visible \cite{page-ban}, and Instagram offers a silent Restrict \cite{ig-restrict}. Personal profiles offer block, a perception-limiting Restricted list \cite{fb-restricted}, and audience-wide comment limits. The separation should be explicit, so that it serves correction and not only concealment. Block is retained for stalking, threats, and cases where visibility itself creates danger.
  \item \textbf{A graduated moderation sequence for conduct.} Pre-posting prompts, return with explanation, removal with notice, temporary exclusion, restricted interaction, and block, so that mild intentions do not collapse to the top rung of the access restrictions.
  \item \textbf{Norm-stating removals.} Public notices that state the rule applied without naming the person, so that moderation leaves a reusable account.
  \item \textbf{Attribution disclosure for attached content.} A visible indication that an advertisement was placed by the platform and not chosen by the author, which may reduce $J_o$ for commercial endorsement without requiring the user to control the advertisement. It may not reduce $A_o$, since association can persist when viewers know the author did not choose the advertisement.
  \item \textbf{Symmetric adjacency controls.} Where advertisers manage adjacency through inventory filters and blocklists, ordinary posters could be offered a corresponding control over what appears beside their posts and Reels, subject to the platform's commercial constraints. A version already exists for some monetized creators \cite{sej-overlay}.
  \item \textbf{Audience-limited publishing without recommendation use.} An option to share with a chosen audience without the material entering the platform's recommendation of content to others. Its feasibility is not shown here.
  \item \textbf{An author-level option to end a Reel's viewing sequence without platform-selected continuation.} An option that returns the viewer to the author's content, or ends playback, instead of advancing to content selected by the recommendation system. Its feasibility and its effect on platform metrics are not shown here.
\end{enumerate}

Each item separates a relation that is currently bundled. None is claimed to be costless, and the simplicity objection applies to all of them.

\section{Conclusion}\label{sec:conclusion}

Interfaces govern in two ways. They permit and prohibit actions, and they decide which meanings can be enacted separately. The second form of power is easy to miss because it appears in no rule. It shows in what a single act carries with it: an advertisement displayed with a post, a block that removes more than it was meant to, participation that also lends support to a system the participant may not endorse, a restriction imposed without the explanation that would make it correction.

The account developed here gives that power a form. An action has an intended set and an effect set. The difference has two parts, surplus, what the action does beyond the intention, and deficit, what the intention required and the action does not do. Mismatch matters when observers attribute it to the user, when the user loses something they value, is made to bear a relation they reject, or fails to obtain a change they rate highly. An action is defective when a cleaner action is feasible and not offered. Applied to blocking, advertising adjacency, presence as support, and recommendation, the test returns different results. It finds the first case provisionally defective on personal profiles, since other Meta surfaces show that the separation is feasible, leaves the second undetermined pending evidence on attribution and availability, finds the third consequential for some users by aversion but undetermined as a defect, and excludes the fourth.

The remedy is not better judgment by users, who are already exercising judgment within a coarse action space. It is a larger action space, with finer separations between what a user does and what the platform makes of it. Where that space is small, one remaining option is to say less, and the loss of that speech is a cost that the design imposes and that may, if the withdrawal hypothesis holds, eventually fall on the platform.
